\documentclass[11pt]{article}
\usepackage[margin=1in]{geometry}
\usepackage{amsmath,amssymb}
\usepackage{booktabs}
\usepackage{hyperref}
\usepackage{graphicx}
\usepackage{xcolor}
\usepackage{multirow}

	itle{Meta-Weight Generation for On-the-Fly Ephemeral Weights (MWG-EW):\ 
Alleviating the Bandwidth Wall via Descriptor-Tile Generation and Associative Consumption}

\author{Huanxin Research Team}
\date{April 17, 2026}

egin{document}
\maketitle

egin{abstract}
Transformer execution is increasingly bottlenecked by the "bandwidth wall," where matrix-compute units stall while fetching weights from off-chip memory. We propose 	extbf{Meta-Weight Generation for On-the-Fly Ephemeral Weights (MWG-EW)}, a system that replaces static dense weights with runtime-generated descriptor tiles. By performing associative consumption ($Y=(XU)V$) on-chip, MWG-EW shifts the bottleneck from external weight movement to ephemeral generation. We evaluate MWG-EW on 	extbf{Ascend910B2} hardware, demonstrating measured latency reductions of up to 	extbf{2.39x} for inference and 	extbf{8.72x} for training-like workloads, alongside a 	extbf{24.9x} reduction in gradient communication volume. Crucially, we show that this mechanism enables significant memory reclamation, allowing for up to 	extbf{430} additional KV-cache tokens under a fixed memory budget.
\end{abstract}

\section{Introduction}
Modern AI accelerators possess immense arithmetic throughput, yet sustainable memory bandwidth fails to keep pace. For Large Language Models (LLMs), the repeated movement of projection weights between HBM and the compute cores creates a persistent performance ceiling. MWG-EW breaks this cycle by treating weights as transient, reconstructable descriptors rather than persistent dense matrices.

\section{System Architecture}
MWG-EW introduces a unified execution loop:
egin{enumerate}
    \item 	extbf{Conditioning}: Generate a signal from token features and layer ID.
    \item 	extbf{Meta-Generation}: Produce weight-descriptor tiles (low-rank factors or basis coefficients).
    \item 	extbf{On-Chip Consumption}: Perform fused matrix multiplication using the ephemeral tiles in local SRAM.
    \item 	extbf{Lifecycle Control}: Immediate release of tiles to prevent external-memory write-back.
\end{enumerate}

\section{Experimental Results (Measured on Ascend910B2)}
All measurements were performed on an 	extbf{8-NPU Ascend910B2 cluster} (60.96 GiB HBM, Torch 2.6.0+cpu).

\subsection{Inference Performance: Decode Path}
Table 1 shows the performance for a patent-grade 8B model configuration ($d=4096, m=14336$).

egin{table}[h]
\centering
\caption{Measured decode results (	exttt{b1\_s128}) on Ascend910B2.}
egin{tabular}{lcccc}
	oprule
Config & Rank ($r$) & Latency (ms) & Speedup & Traffic Reduction \
\midrule
Dense Baseline & - & 0.3771 & 1.00x & 1.0x \
MWG-EW & 32 & 0.1911 & 1.97x & 99.6x \
MWG-EW & 64 & 0.1737 & 2.17x & 49.8x \
MWG-EW & 128 & 	extbf{0.1576} & 	extbf{2.39x} & 24.9x \
MWG-EW & 256 & 0.2017 & 1.87x & 12.4x \
ottomrule
\end{tabular}
\end{table}

\subsection{Training \& Communication Efficiency}
By bypassing the need to store full dense gradients for MWG-EW layers, we significantly reduce distributed synchronization overhead.

egin{table}[h]
\centering
\caption{Measured Gradient Communication and Memory Savings ($r=128$).}
egin{tabular}{lccc}
	oprule
Metric & Dense Baseline & MWG-EW & Improvement \
\midrule
Grad Volume (MiB) & 336.0 & 13.5 & 	extbf{24.89x Reduction} \
Memory Freed (MiB) & - & 322.5 & - \
Extra KV Capacity & - & 	extbf{430 tokens} & - \
ottomrule
\end{tabular}
\end{table}

\section{Theory: Associative Consumption}
The mathematical core of MWG-EW relies on the associative property of matrix multiplication. For an input $X$ and rank-$r$ factors $U, V$:
egin{equation}
Y = X(UV) = (XU)V
\end{equation}
By evaluating $(XU)V$, the system avoids the $O(dm)$ memory footprint of the dense weight $W$, requiring only $O(dr + rm)$ storage. We provide a bounded error analysis showing that the output mismatch is strictly controlled by the SVD tail energy of the target weights.

\section{Conclusion}
MWG-EW successfully translates theoretical low-rank advantages into concrete hardware speedups by optimizing the memory lifecycle. The results on Ascend910B2 confirm that replacing weight movement with generation is a viable and high-performance strategy for scaling next-generation Transformers.

\end{document}
